\par\Needspace{6\baselineskip}\medskip
\SourceLeaf{158}\pdfbookmark[2]{113.}{article-113}
\begin{center}113.\end{center}
Damit die Formeln nicht allzu umfangreich werden, verwenden wir die folgenden Abkürzungen. Zunächst bezeichnen wir die Größe
\[
\operatorname{tang}\beta\sin(\alpha''-\alpha')+\operatorname{tang}\beta'\sin(\alpha-\alpha'')+\operatorname{tang}\beta''\sin(\alpha'-\alpha)
\]
mit $(0.1.2)$. Ersetzt man in diesem Ausdruck Länge und Breite eines beliebigen geozentrischen Ortes durch Länge und Breite eines beliebigen der drei heliozentrischen Orte der Erde, so ersetzen wir im Zeichen $(0.1.2)$ die dem ersten Ort entsprechende Zahl durch die römische Zahl für den eingesetzten Ort. So bezeichnet etwa das Zeichen $(0.1.\mathrm{I})$ die Größe
\[
\operatorname{tang}\beta\sin(L'-\alpha')+\operatorname{tang}\beta'\sin(\alpha-L')+\operatorname{tang}B'\sin(\alpha'-\alpha)
\]
und das Zeichen $(0.\mathrm{O}.2)$ die Größe
\[
\operatorname{tang}\beta\sin(\alpha''-L)+\operatorname{tang}B\sin(\alpha-\alpha'')+\operatorname{tang}\beta''\sin(L-\alpha).
\]
Ebenso ändern wir das Zeichen, wenn im ersten Ausdruck \emph{zwei} geozentrische Paare aus Länge und Breite durch zwei beliebige heliozentrische Paare der Erde ersetzt werden. Werden zwei Paare aus Länge und Breite, die in denselben Ausdruck eingehen, lediglich miteinander vertauscht, so muss man auch im Zeichen die entsprechenden Zahlen vertauschen. Der Wert selbst ändert sich dabei nicht, sondern wird nur aus einem positiven zu einem negativen oder aus einem negativen zu einem positiven. So gilt etwa $(0.1.2)=-(0.2.1)=(1.2.0)=-(1.0.2)=(2.0.1)=-(2.1.0)$. Alle auf diese Weise entstehenden Größen lassen sich daher auf die folgenden 19 zurückführen
\[
\begin{aligned}
&(0.1.2)\\
&(0.1.\mathrm{O}),\ (0.1.\mathrm{I}),\ (0.1.\mathrm{II}),\\
&(0.\mathrm{O}.2),\ (0.\mathrm{I}.2),\ (0.\mathrm{II}.2),\\
&(\mathrm{O}.1.2),\ (\mathrm{I}.1.2),\ (\mathrm{II}.1.2)\\
&(0.\mathrm{O}.\mathrm{I}),\ (0.\mathrm{O}.\mathrm{II}),\ (0.\mathrm{I}.\mathrm{II}),\\
&(1.\mathrm{O}.\mathrm{I}),\ (1.\mathrm{O}.\mathrm{II}),\ (1.\mathrm{I}.\mathrm{II}),\\
&(2.\mathrm{O}.\mathrm{I}),\ (2.\mathrm{O}.\mathrm{II}),\ (2.\mathrm{I}.\mathrm{II}),
\end{aligned}
\]
Zu ihnen kommt als zwanzigste $(\mathrm{O}.\mathrm{I}.\mathrm{II})$ hinzu.
\par
Es lässt sich leicht zeigen, dass jeder dieser Ausdrücke, multipliziert mit dem Produkt der Kosinus der drei in ihn eingehenden Breiten, gleich dem sechsfachen Volumen der Pyramide ist, deren Spitze in der Sonne liegt und deren Grundfläche das Dreieck aus den drei Punkten der Himmelskugel ist, die den in den Ausdruck eingehenden Orten entsprechen, wenn man den Radius der Kugel $=1$ setzt. Liegen also diese drei \SourceLeaf{159}Orte auf demselben Großkreis, so muss der Wert des Ausdrucks $=0$ sein. Dies trifft für die drei heliozentrischen Orte der Erde stets zu, wenn wir die Parallaxen und die durch Störungen verursachten Breiten der Erde außer Betracht lassen, das heißt, wenn wir die Erde in der Ebene der Ekliptik selbst ansetzen. Unter dieser Annahme gilt daher stets $(\mathrm{O}.\mathrm{I}.\mathrm{II})=0$. Diese Gleichung ist eine Identität, wenn die Ekliptik selbst als dritte Ebene gewählt wurde. Verschwinden ferner $B$ und $B'$ und gilt $B''=0$, so werden alle diese Ausdrücke bis auf den ersten viel einfacher. Vom zweiten bis zum zehnten bestehen sie jeweils aus zwei Gliedern, vom elften bis zum neunzehnten jeweils aus einem einzigen Glied.
